% ============================================================
%  THE ORTYX PROTOCOL
%  Appendix B: Salience Fields and Jitter Geometry
% ============================================================

\chapter{Salience Fields and Jitter Geometry}
\label{app:salience}

\noindent
This appendix develops the formal mathematics of the
Marine Salience Algorithm and the salience manifold
structure. These definitions provide the operational
grounding for the \LevelI{} claims about the Marine
metric that survive the decomposition of the Phoenix
corpus.

\section{Signal Representation and Peak Detection}
\label{app-b:peaks}

Let a signal be represented as a continuous waveform
$s : \mathbb{R} \to \mathbb{R}$. A peak detection
procedure identifies a sequence of local extrema:
\[
  \{(t_i, a_i)\}_{i=1}^{N},
\]
where $t_i$ denotes the peak time and $a_i = s(t_i)$
the peak amplitude, for $t_1 < t_2 < \cdots < t_N$.

\section{Jitter Metrics}
\label{app-b:jitter}

Define the instantaneous period at peak $i$:
\[
  T_i = t_i - t_{i-1}, \quad i \ge 2.
\]
Define exponential moving averages with smoothing
parameters $\alpha, \beta \in (0,1)$:
\[
  \bar{T}_i = (1-\alpha)\bar{T}_{i-1} + \alpha T_i,
  \qquad
  \bar{A}_i = (1-\beta)\bar{A}_{i-1} + \beta a_i.
\]
The period jitter and amplitude jitter at peak $i$ are:
\[
  \Jp(i) = |T_i - \bar{T}_i|, \qquad
  \Ja(i) = |a_i - \bar{A}_i|.
\]
These measure the deviation of the current peak from
the recent local trend in period and amplitude respectively.

\begin{definition}
A signal $s(t)$ is \emph{locally $\delta$-coherent}
at peak $i$ if $\Jp(i) + \Ja(i) < \delta$.
\end{definition}

\section{The Salience Functional}
\label{app-b:salience-functional}

Let $E_i$ denote the local signal energy at peak $i$
(e.g., a windowed RMS value around $t_i$) and $H_i$
a harmonic alignment score (measuring consistency with
a harmonic series). The Marine salience functional is:
\[
  \SalFunc(i) = w_E E_i + w_H H_i
  + w_J \cdot \frac{1}{1 + \Jp(i) + \Ja(i)},
\]
where $w_E, w_H, w_J > 0$ are weighting parameters
with $w_E + w_H + w_J = 1$ (by normalization convention).

\begin{proposition}
The salience functional $\SalFunc(i)$ is bounded:
$\SalFunc(i) \in [0, w_E \|E\|_\infty + w_H + w_J]$.
In the limit of perfect temporal regularity
($\Jp(i) + \Ja(i) \to 0$), the jitter term achieves
its maximum value $w_J$.
\end{proposition}

\section{The Salience Manifold}
\label{app-b:salience-manifold}

For a threshold $\lambda > 0$, define the coherent
salience fiber:
\[
  \mathcal{F}_\lambda = \{i : \SalFunc(i) > \lambda\}.
\]
The salience manifold is the union over all thresholds:
\[
  \SalMan = \bigcup_{\lambda > 0} \mathcal{F}_\lambda.
\]
In practice, $\SalMan$ is the set of all peak indices
that exceed some minimum salience threshold $\lambda_{\min}$.
The Marine filter passes only events in $\SalMan$
to downstream processing.

\section{Computational Complexity}
\label{app-b:complexity}

The Marine algorithm as described requires, per peak:
two multiplications and additions (EMA updates),
two absolute values (jitter computation), one reciprocal
(coherence term), and three multiply-accumulate operations
(weighted sum). This is $O(1)$ operations per peak,
with no dependence on the signal length or any global
state beyond the running averages $\bar{T}_i$ and
$\bar{A}_i$.

\begin{proposition}
The Marine Salience Algorithm achieves $O(1)$
time and $O(1)$ space per peak (excluding the output
event sequence), compared to $O(N \log N)$ time for
FFT-based spectral analysis with frame length $N$.
\end{proposition}

\section{Jitter and Accessibility Entropy}
\label{app-b:jitter-entropy}

The RSVP reinterpretation of the jitter metric
(Chapter~17) proposes the structural relationship:
\[
  \SalFunc(i) \propto \frac{1}{\AccEnt(x_{t_i})},
\]
where $\AccEnt(x_t) = \log|\AdmTraj(x_t)|$ is the
accessibility entropy of the signal state at time $t_i$.

This proportionality is a research hypothesis rather
than a derivation. For it to become a derivation,
one would need to: (1)~specify a state space $X$
for audio signals in which trajectories are well-defined,
(2)~define a measure on $\AdmTraj(x_t)$ consistent
with the signal's physical dynamics, and (3)~demonstrate
that low-jitter peaks correspond to low-$\AccEnt$
states under this measure. These are open research
problems that constitute the \LevelI{} prediction
generated by the RSVP reinterpretation.
